% =============================================================
%                ESTADO ARQUITETURAL DO MODELO
% =============================================================
\section{Estado arquitetural do modelo}

O estado funcional do coprocessador contém:

\begin{itemize}
    \item 32 registradores vetoriais físicos;
    \item o valor corrente de \texttt{vl};
    \item o valor corrente de \texttt{vtype}, composto por SEW e LMUL;
    \item a memória usada pelos testes de load e store.
\end{itemize}

\subsection{Representação física dos registradores}

Cada registrador vetorial físico é representado como uma sequência de
16 bytes. Essa representação permite reinterpretar o mesmo estado físico
sob diferentes valores de SEW.

Por exemplo, os mesmos 16 bytes correspondem a:

\begin{itemize}
    \item 16 elementos de 8 bits;
    \item 8 elementos de 16 bits;
    \item 4 elementos de 32 bits.
\end{itemize}

A ordem dos bytes dentro de um elemento é little-endian.

\subsection{Grupos de registradores}

Quando LMUL é maior que 1, um registrador lógico ocupa um grupo
consecutivo de registradores físicos. O tamanho do grupo é igual a LMUL.

O índice inicial do grupo deve ser múltiplo de LMUL. O grupo também deve
permanecer dentro dos 32 registradores físicos.

Exemplos:

\begin{itemize}
    \item LMUL igual a 2: índices-base válidos são 0, 2, 4, \ldots, 30;
    \item LMUL igual a 4: índices-base válidos são 0, 4, 8, \ldots, 28;
    \item LMUL igual a 8: índices-base válidos são 0, 8, 16 e 24.
\end{itemize}

A leitura e escrita de elementos lógicos concatenam, em ordem crescente,
o conteúdo dos registradores físicos que compõem o grupo.

% =============================================================
%                  ESPECIFICAÇÃO EXECUTÁVEL
% =============================================================
\section{Especificação executável}

\subsection{Programa}

O programa funcional é representado como uma lista ordenada de pares
compostos por mnemônico e operandos nomeados.

Exemplo conceitual:

\begin{verbatim}
[
    ("vsetvli", {"vl": 4, "sew": 32, "lmul": 1}),
    ("vadd.vv", {"vd": 3, "vs1": 1, "vs2": 2})
]
\end{verbatim}

A classe \texttt{Instruction} valida o mnemônico e os operandos de cada
instrução antes da execução. O executor processa a lista em ordem,
aplicando cada instrução ao estado do modelo.

\subsection{Execução sequencial}

O executor atual é funcional e sequencial:

\begin{enumerate}
    \item valida ou normaliza a instrução;
    \item identifica o handler correspondente;
    \item executa a operação sobre o modelo;
    \item registra os efeitos observáveis no trace.
\end{enumerate}

Não há modelagem de ciclo, pipeline, issue múltiplo, hazards,
interrupções, flush ou concorrência de memória.

\subsection{Trace}

O trace registra uma entrada por instrução executada. Os campos atuais
incluem:

\begin{itemize}
    \item índice da instrução;
    \item mnemônico;
    \item operandos;
    \item valor de \texttt{vl} durante a execução;
    \item registradores de destino observados;
    \item escritas de memória observadas, quando aplicável.
\end{itemize}

O trace pode ser serializado em JSON. Valores incluídos no JSON devem
ser tipos serializáveis, como inteiros, strings, listas e dicionários.
Valores de memória não devem ser armazenados como objetos \texttt{bytes}
sem uma conversão explícita.

% =============================================================
%                       STRIP-MINING
% =============================================================
\section{Strip-mining e tamanho residual}

O suporte atual a strip-mining gera uma sequência de instruções para
processar um vetor lógico maior que VLMAX.

Para um número total de elementos |$N$|, cada iteração escolhe:

|$$
vl_{\text{iter}} = \min(N_{\text{remaining}},VLMAX).
$$|

Após cada iteração:

|$$
offset_{\text{elements}} \leftarrow
offset_{\text{elements}} + vl_{\text{iter}},
$$|

|$$
N_{\text{remaining}} \leftarrow
N_{\text{remaining}} - vl_{\text{iter}}.
$$|

Um kernel de soma contígua executa, em cada iteração:

\begin{enumerate}
    \item configuração de \texttt{vl};
    \item load do trecho corrente do vetor A;
    \item load do trecho corrente do vetor B;
    \item soma dos registradores vetoriais;
    \item store do trecho corrente do resultado.
\end{enumerate}

Para VLEN igual a 128, SEW igual a 32 e LMUL igual a 1,
|$VLMAX=4$|. Um vetor de 10 elementos é processado com comprimentos:

|$$
4,\ 4,\ 2.
$$|

O avanço de endereço entre iterações é calculado em bytes:

|$$
offset_{\text{bytes}} =
offset_{\text{elements}} \times element\_bytes.
$$|

A implementação inicial de strip-mining foi escrita para SEW igual a
32 e LMUL igual a 1. Sua generalização para SEW e LMUL variáveis ainda
é uma pendência.